\chapter{Justification}\label{cap:justification}

Although numerous works have attempted to define a function for adult neurogenesis, the role played by the immature cells it
generates within the circuit they integrate into remains poorly understood~\parencite{aimoneComputational2016, aimoneModeling2011,
berdugo-vegaSharpening2023, faresNeurogenesis2019}.

A number of studies on this topic have focused primarily on the effects of neurogenesis in the DG itself and on the pattern
separation it performs~\parencite{berdugo-vegaSharpening2023, kimAdult2024, wangEffect2024}, neglecting its impact on connected
regions such as the CA3 area. This work therefore aims to fill this gap by also examining how iGCs affect the sparsity and the
pattern separation of CA3, and how the backprojection from this area to the DG influences the
system~\parencite{myersPattern2011}, with the goal of producing a more holistic and biologically plausible model.

Extending the model beyond the DG is also what makes the competing accounts separable. The circuit mechanisms proposed for
adult-born cells, restraining activity in the surrounding DG and differentially engaging CA3 interneurons among them, predict the
same behavioral outcome and can only be told apart in a model that contains both regions and the interneurons connecting
them~\parencite{aimoneComputational2016, wangEffect2024}. A model without synaptic plasticity adds a further distinction that the
experimental preparations cannot easily make, since any effect it reproduces is attributable to the circuit rather than to
learning.

This study is therefore justified by filling a gap in the existing literature, carrying the modeling of adult neurogenesis beyond
the DG to the circuit through which its effects reach CA3.
